\SourceLeaf{271}%
die Anwendung auf eine bestimmte Zone anzusehen sind. Einige Umformungen, die dabei helfen, verdienen daher hier einen Platz. Eine Umformung ist ohnehin \textit{notwendig}, wenn man von einer bestimmten Normalbreite auf der Kugel anstelle derjenigen auf dem Ellipsoid ausgeht, also von einem gegebenen Wert für \(Q\), und daraus die übrigen Größen berechnen will.

Führt man drei Hilfswinkel \(\varphi,\zeta,\eta\) so ein, dass
\[
\sin\varphi=e \quad .\ .\ .\ .\ .\ .\ .\ .\ .\ .\ .\ .\ .\ .\ .\quad(13)
\]
\[
\operatorname{tang}\zeta=\operatorname{tang}\varphi\cos P^2 \quad .\ .\ .\ .\ .\ .\ .\ .\ .\ .\quad(14)
\]
\[
\operatorname{tang}\eta=\sin\zeta\operatorname{tang}P \quad .\ .\ .\ .\ .\ .\ .\ .\ .\ .\ .\quad(15)
\]
so erhält man
\[
\alpha=\frac{1}{\cos\zeta} \quad .\ .\ .\ .\ .\ .\ .\ .\ .\ .\ .\ .\ .\ .\quad(16)
\]
\[
\sin Q=\cos\zeta\sin P \quad .\ .\ .\ .\ .\ .\ .\ .\ .\ .\ .\ .\quad(17)
\]
\[
\cos\eta\cos Q=\cos P \quad .\ .\ .\ .\ .\ .\ .\ .\ .\ .\ .\ .\ .\quad(18)
\]
\[
\sin\eta=\operatorname{tang}\zeta\operatorname{tang}Q \quad .\ .\ .\ .\ .\ .\ .\ .\ .\ .\quad(19)
\]
\[
\operatorname{tang}\tfrac{1}{2}(P-Q)=\operatorname{tang}\tfrac{1}{2}\zeta\mathord{.}\operatorname{tang}\tfrac{1}{2}\eta \quad .\ .\ .\ .\ .\ .\ .\ .\ .\quad(20)
\]
\[
\sin(2\zeta-\varphi)=e\cos2Q \quad .\ .\ .\ .\ .\ .\ .\ .\ .\ .\ .\ .\quad(21)
\]

Gleichung 18 folgt leicht aus der Verbindung von 15 und 17. Ebenso ergibt sich 19 aus 15, 17 und 18, dann 20 aus 17, 18 und 19 und schließlich 21 aus 14 und 17.

Am genauesten rechnet man, wenn man bei gegebenem \(P\) die Formeln 14, 15 und 20 verwendet, um der Reihe nach \(\zeta,\eta,Q\) zu bestimmen. Ist dagegen \(Q\) gegeben, so leitet man aus den Gleichungen 21, 19 und 20 die Werte von \(\zeta,\eta,P\) ab. Zur Kontrolle kann man dann noch eine oder mehrere der übrigen Gleichungen verwenden. Führt man einen vierten Hilfswinkel \(\theta\) nach der Formel
\[
\sin\theta=e\sin P \quad .\ .\ .\ .\ .\ .\ .\ .\ .\ .\ .\ .\ .\ .\quad(22)
\]
ein, so erhält man
\[
\cos\varphi=\cos\zeta\cos\eta\cos\theta \quad .\ .\ .\ .\ .\ .\ .\ .\ .\ .\quad(23)
\]
\Needspace{9\baselineskip}
und die Formeln 9 und 12 nehmen die folgende Gestalt an.
\[
A=\frac{a\cos P}{\alpha\cos\theta\cos Q}=\frac{a\cos\eta}{\alpha\cos\theta}=\frac{a\cos\varphi}{\cos\theta^2}=\frac{a\cos\varphi}{1-ee\sin P^2}
\]
\[
k=\frac{\operatorname{tang}(45^0+\frac{1}{2}P)^\alpha}{\operatorname{tang}(45^0+\frac{1}{2}Q)\operatorname{tang}(45^0+\frac{1}{2}\theta)^{\alpha e}}
\]
